\documentclass[11pt,a4paper]{article}
\usepackage[T1]{fontenc}
\usepackage[utf8]{inputenc}
\usepackage{lmodern,microtype}
\usepackage[margin=27mm,headheight=14pt]{geometry}
\usepackage{amsmath,amssymb,amsthm,mathtools,mathrsfs}
\usepackage{booktabs,array,longtable,enumitem}
\usepackage{xcolor}
\usepackage{xurl}
\usepackage{hyperref}
\usepackage{fancyhdr}
\hypersetup{colorlinks=true,linkcolor=blue!45!black,citecolor=blue!45!black,urlcolor=blue!45!black,
 pdftitle={Arithmetic Transseries Beyond an Accumulation Cut},
 pdfauthor={Research manuscript prepared for Vladimir Reshetnikov}}
\urlstyle{same}
\pagestyle{fancy}\fancyhf{}
\fancyhead[L]{\small Arithmetic transseries beyond an accumulation cut}
\fancyhead[R]{\small 29 September 2026}
\fancyfoot[C]{\thepage}
\setlength{\parskip}{4pt}
\setlength{\emergencystretch}{1.5em}
\setlength{\parindent}{14pt}
\setlist[enumerate]{itemsep=3pt,topsep=4pt}
\setcounter{tocdepth}{2}
\numberwithin{equation}{section}
\newtheorem{theorem}{Theorem}[section]
\newtheorem{proposition}[theorem]{Proposition}
\newtheorem{lemma}[theorem]{Lemma}
\newtheorem{corollary}[theorem]{Corollary}
\theoremstyle{definition}
\newtheorem{definition}[theorem]{Definition}
\newtheorem{example}[theorem]{Example}
\newtheorem{question}{Research question}
\theoremstyle{remark}
\newtheorem{remark}[theorem]{Remark}
% ed. (2026-09-29): unnumbered environment for ProveIt editorial notes; it
% leaves every numbered statement of the delivered article unchanged.
\newtheorem*{ednote}{Editorial note (ProveIt, 2026-09-29)}
\newcommand{\N}{\mathbb N}
\newcommand{\C}{\mathbb C}
\newcommand{\Q}{\mathbb Q}
\newcommand{\R}{\mathbb R}
\newcommand{\A}{\mathcal A}
\newcommand{\B}{\mathcal B}
\newcommand{\HH}{\mathscr H}
\newcommand{\ee}{\mathrm e}
\newcommand{\lcm}{\operatorname{lcm}}
\newcommand{\ind}{\mathbf 1}
\newcommand{\ord}{\operatorname{ord}}
\newcommand{\Res}{\operatorname{Res}}
\newcommand{\bigO}{\mathcal O}
\newcommand{\eps}{\varepsilon}
\newcommand{\abs}[1]{\left|#1\right|}
\newcommand{\norm}[1]{\left\|#1\right\|}
\newcommand{\coeff}[2]{[#1^{#2}]}
\newcommand{\snapshot}{\texttt{db68f0853c3c69cd930caedab6f6ad9addb11eaf}}
\begin{document}
% ed. (2026-09-29): no hyperref page anchor on the title page, whose page
% number 1 recurs after it (removes a duplicate-destination warning).
\hypersetup{pageanchor=false}
\begin{titlepage}
\vspace*{12mm}
\noindent\textsc{Transseries \; / \; Arithmetic asymptotics \; / \; Nonlinear inversion}
\vspace{15mm}
\begin{center}
{\LARGE\bfseries Arithmetic Transseries\\[4mm]
Beyond an Accumulation Cut\par}
\vspace{8mm}
{\Large Divisibility, summable inversion, and\\[2mm]
curvature-lifted resonances\par}
\vspace{14mm}
{\large Research manuscript prepared for Vladimir Reshetnikov\par}
\vspace{3mm}
29 September 2026
\end{center}
\vspace{12mm}
\begin{abstract}
Divisor-counting sequences have corrections with arithmetic coefficients
$\ind_{d\mid n}$ and exponential actions $a(1-1/d)$ accumulating at $a$.
This combination defeats finite-grid truncation without defeating Hahn
well-ordering. We construct an arithmetic Hahn ring and a compatible scale
of weighted Banach algebras that retain the indicators exactly. An lcm
convolution handles arithmetic multiplication, while an action--label weight
makes the entire accumulating block summable on expanding families of
analytic sheets. For the elementary input block, membership has the exact
criterion $\sum_d |c_d|\ee^{-ad}<\infty$.

We prove a convergent, all-action, fixed-label inverse theorem for
$\ee^{ax}(1+E_n(x))/x$, with explicit coefficients and a geometric remainder.
It applies to necklaces and primitive necklaces on their exact ranges.
A second result identifies the first inverse sector beyond the primitive
action cut: it has action $a$, but its leading pure-exponential amplitude
cancels, leaving a curvature-induced factor of order $X^{-1}$.
An exact coefficient transformation explains this phenomenon at every
homogeneous order and gives its extension to cores $\ee^{ax}x^{-b}$.

The construction answers a specified algebra-and-summation direction in the
ProveIt companion volume. It is not a canonical interpolation of an
arithmetic sequence, a field containing nontrivial indicators, or a general
resurgence theorem. Full mathematical proofs are supplied; 964 exact finite
assertions and 48 high-precision error checks provide supplementary audits.
These results are not Lean-verified, and publication-level priority has not
been established.
\end{abstract}
\vfill
\noindent\small\textbf{Status.} AI-assisted research manuscript; unrefereed.
Classical ingredients and proposed extensions are distinguished in
Section~\ref{sec:scope}. The repository snapshot, reproducibility instructions,
and formalization boundaries are recorded explicitly.
\end{titlepage}
% ed. (2026-09-29): page anchors resume after the title page.
\hypersetup{pageanchor=true}
{\small\setlength{\parskip}{0pt}\tableofcontents}
\newpage

\section{The research target and the scope of the answer}\label{sec:scope}

\subsection{A specific question in the repository}

The current ProveIt transseries project separates its generic Lean modules
from the Fabius-function development. Its companion volume,
\emph{Combinatorial Transseries and Their Inverses}, contains a chapter on
necklaces, Lyndon words, and irreducible polynomials. The subsequent
``Further research directions'' section explicitly asks for an arithmetic
transseries algebra that preserves divisibility indicators, permits an
accumulation of actions at $\log q$, and specifies summable products and
inverse operations \cite{repo}. This is the target of this article.

The source consulted was the file named in Appendix~\ref{app:provenance},
at the pinned commit \texttt{db68f0853c3c} (full hash in the appendix).
In that source the question occurs near
line 4971. Its benchmark is Proposition
\path{t2:prop:necklace-cutoff}; the relevant chapter has label
\texttt{ct:ch:necklaces}. The source already provides finite divisor cutoffs,
finite-cutoff trigonometric interpolations, certain fixed-radical analytic
sheets, and integer threshold methods. None of those constructions is
claimed here as a new result.

The repository also contains an unmerged article on action accumulation
and nonlinear inversion using Laplace measures. Its own catalogue entry
states that it does not preserve the divisibility indicators and therefore
does not provide the arithmetic algebra requested by the companion
\cite{repoindex}. The present construction addresses that missing
arithmetic component rather than repeating the measure-valued theory.

% ed. (2026-09-29): records the second repository question this article
% answers without citing it, and the relation to the action-measure article,
% which the article knew only from a catalogue line.
\begin{ednote}\label{ed:siblings}
\emph{A second question answered.} The same construction answers the
problem ``Arithmetic action accumulation rather than a discrete Borel
lattice'' of the inverse-harmonic package
(\path{Analysis/Transseries/docs/series-and-transseries/Inverse_Harmonic_Stokes_Transport/inverse_harmonic_transseries.tex}),
which asks for a summable inverse calculus for divisor-dependent sectors
whose actions accumulate at a finite point, with a specified topology or
summability condition, retained divisibility indicators, and justified
products and inverse operations. Theorems~\ref{thm:banach} and
\ref{thm:admission} supply the topology and the products,
Corollary~\ref{cor:funcalc} the scalar operations, and
Theorems~\ref{thm:inverse} and \ref{thm:range} the inverses, within the
limits stated in this article: inverses on each fixed arithmetic sheet, no
canonical global interpolation, and no resurgence, Borel summation or
Stokes automorphism. This article does not cite that problem.
\emph{The action-measure article.} The ``catalogue entry'' cited above is
the repository's description of
\path{Analysis/Transseries/docs/series-and-transseries/Action_Accumulation_Nonlinear_Inversion/},
not a statement of that article itself, and this article did not read it.
That article gives an analytic inversion calculus for exponentially
weighted action measures with the explicit inverse
$-\sum_{n\ge1}s^{n-1}\mu^{*n}/n!$ and no action-gap hypothesis, a
Poisson--Bessel resolution of $\sum_j j^{-2}\ee^{-x/j}$, and Hardy-field
and o-minimal obstructions. It is label-free and analytic; this article is
fixed-label and arithmetic. No theorem of either is a theorem of the
other, and neither supersedes the other: they approach the same
accumulation frontier from opposite sides and form one unit for the
deferred merge into the companion's necklace chapter.
\end{ednote}

\subsection{What is established}

There are three main results, together with their necessary limitations.

\paragraph{Arithmetic summation across the cut.}
Theorems~\ref{thm:banach} and \ref{thm:admission} provide an explicit
weighted Banach algebra whose coefficients multiply by lcm convolution.
The whole infinite family of primitive actions is retained. Scalar analytic
functional calculus, including logarithms and multiplicative inverses near
one, is norm-convergent. Evaluation gives exact arithmetic data on the
diagonal $x=n$ and holomorphic functions on label-dependent half-planes.

\paragraph{Convergent inverse transport.}
Theorems~\ref{thm:inverse} and \ref{thm:range} give an explicit inverse
series around the exact core $\ee^{aX}/X$. The series is absolutely
summable even after it is expanded into individual action tuples. A
uniform bound controls its tail. The arithmetic label is fixed throughout
an inverse calculation; this is an essential part of the statement.

\paragraph{Curvature-lifted resonance.}
Theorems~\ref{thm:firstcut} and \ref{thm:resonance} identify a cancellation
at the first action beyond the primitive block and explain it at every
homogeneous order. For the more general core $\ee^{ax}x^{-b}$,
Corollary~\ref{cor:bcore} gives the first surviving coefficient at each
pure-exponential resonance.

The first two results supply a constructive answer to the formal,
analytic-functional-calculus, and fixed-label reversion portions of the
repository question. They deliberately do not postulate a fractional
shift of the arithmetic label. Section~\ref{sec:obstructions} proves why
ordinary differential-field and compact-profinite interpretations cannot
simply be imposed on these data.

\subsection{Classical ingredients and proposed extensions}

The multiplication
$\sum_{\lcm(i,j)=l}c_i b_j$ is the classical lcm convolution. Its
relationship with the divisor-zeta transform is recorded, for example, in
T\'oth's survey, Section 4.4 \cite{toth}. It is not introduced here as a
new arithmetic operation. Hahn supports, well-ordering arguments, and
Lagrange inversion are also classical ingredients; Edgar's account gives
background on transseries and their usual differential and compositional
operations \cite{edgar}.

The proposed contributions are the combined action--label completion,
its exact elementary-block admission criterion and faithful sheet
realization, the absolute inverse transport through the accumulating
arithmetic support, and the all-order curvature calculation and its
post-cut consequence. Each has a proof below. A targeted literature check
also found a 2026 preprint on analytic properties of necklace polynomials;
it studies the alphabet-variable polynomial and its analytic inequalities,
rather than the length-variable construction considered here
\cite{chebolu}. This comparison is not an exhaustive novelty search. The
results should be regarded as proposed research extensions, not as a
verified claim of first discovery.

\section{Divisor sectors and the two variables}\label{sec:data}

Throughout, $\N=\{1,2,\ldots\}$ and $a>0$. For integer alphabet size
$q\ge2$ we specialize to $a=\log q$. Set
\begin{equation}
 e_d(n)=\ind_{d\mid n},\qquad
 \lambda_d=a\left(1-\frac1d\right),\qquad d\ge2.
 \label{eq:actions}
\end{equation}
The letter $n$ always denotes an arithmetic label. The variable $x$ is a
real or complex analytic coordinate. They can be set equal, but they are
not identified while differentiating or reverting a function.

\subsection{The exact combinatorial input}

The number of necklaces and primitive necklaces of length $n$ are,
respectively,
\begin{equation}
 N_q(n)=\frac1n\sum_{d\mid n}\varphi(d)q^{n/d},\qquad
 L_q(n)=\frac1n\sum_{d\mid n}\mu(d)q^{n/d}.
 \label{eq:counts}
\end{equation}
For completeness, the first formula follows by averaging the fixed words
of the cyclic rotations. A rotation through $r$ positions fixes
$q^{\gcd(n,r)}$ words, and grouping the rotations by
$d=n/\gcd(n,r)$ gives $\varphi(d)$ copies. For the second formula, let
$P_q(n)$ count words of least period $n$. Every word has a unique least
period dividing $n$, so $q^n=\sum_{d\mid n}P_q(d)$. M\"obius inversion
then gives the numerator of the second expression, and every primitive
orbit has $n$ words. These are the same exact inputs used in the
repository chapter \cite{repo}.

Given coefficients $c_d$, define
\begin{equation}
 E_n(x)=\sum_{\substack{d\mid n\\d\ge2}}c_d\ee^{-\lambda_d x},
 \qquad f_n(x)=\frac{\ee^{ax}}x\bigl(1+E_n(x)\bigr).
 \label{eq:sheets}
\end{equation}
Each $E_n$ is an entire function because it is a finite sum. Choosing
$c_d=\varphi(d)$ or $c_d=\mu(d)$ gives $f_n(n)=N_q(n)$ or $L_q(n)$.
The family $f_n$ is not asserted to be the restriction of a single
interpolation of the counting sequence.

Two bounds used repeatedly below are
\begin{equation}
 \sum_{d\mid n}\varphi(d)=n,
 \qquad \sum_{d\mid n}|\mu(d)|\le n.
 \label{eq:arithbounds}
\end{equation}
The first identity counts the integers in $\{1,\ldots,n\}$ by the order
of the corresponding element of the cyclic group. The second follows
because each summand is at most one and $n$ has at most $n$ divisors.
Consequently, for either of the two counting inputs and real $X>0$,
\begin{equation}
 A_n(X):=\sum_{\substack{d\mid n\\d\ge2}}|c_d|\ee^{-\lambda_dX}
 \le n\ee^{-aX/2}.
 \label{eq:An}
\end{equation}

\subsection{Why the accumulation point matters}

The actions obey
\begin{equation}
 \frac a2=\lambda_2<\lambda_3<\cdots<a,
 \qquad \lambda_d\longrightarrow a.
 \label{eq:cut}
\end{equation}
A finite additive grid with strictly positive generators has only
finitely many elements below any fixed bound. The set in
\eqref{eq:cut} does not. Moreover $2\lambda_2=a$, so infinitely many
primitive sectors precede the first quadratic sector.

This is not a failure of well-ordering: an increasing sequence can be
well-ordered while accumulating at its supremum. The distinction between
\emph{finitely many decompositions of one action} and \emph{finitely many
actions below a bound} will be central. The former is true here; the
latter is false.

For a finite divisor cutoff $D$, \eqref{eq:arithbounds} gives
\begin{equation}
 \abs{E_n(n)-\sum_{2\le d\le D}c_d e_d(n)\ee^{-\lambda_dn}}
 \le n\ee^{-\lambda_{D+1}n}.
 \label{eq:cutoff}
\end{equation}
This recovers the companion's benchmark without changing its scope.
The goal is now to sum the full primitive block and transport it through
nonlinear operations, not to pretend that \eqref{eq:cutoff} resolves every
bounded action with a fixed finite list.

\section{The arithmetic Hahn ring}\label{sec:formal}

\subsection{The action monoid}

Let
\begin{equation}
 S=\left\{\lambda_{d_1}+\cdots+\lambda_{d_m}:
 m\ge0,\ d_i\ge2\right\},
 \label{eq:S}
\end{equation}
where the empty sum is zero. We order $S$ by its usual real order.

\begin{lemma}[Finite action decompositions]\label{lem:support}
The set $S$ is well-ordered. Every $s\in S$ has only finitely many ordered
representations by generators $\lambda_d$, including the length of the
representation. Every fixed action therefore has finitely many
decompositions into any prescribed finite number of elements of $S$.
\end{lemma}
\begin{proof}
Every representation of an action at most $B$ has length at most
$\lfloor 2B/a\rfloor$. For a fixed length $m$, use the elementary
well-quasi-order property of $\N^m$: every infinite sequence of tuples
has an earlier tuple that is componentwise at most a later one. One proof
is to successively extract an infinite nondecreasing subsequence in each
coordinate. For a single coordinate, either a value occurs infinitely
often or one can successively choose strictly increasing values.

An infinite strictly decreasing sequence of actions in $S$ is bounded
above by its first term. After selecting a subsequence, its chosen
representations have a common length $m$. A componentwise comparable
pair of tuples contradicts the strict decrease, because $d\mapsto
\lambda_d$ is strictly increasing. Thus there is no such decreasing
sequence. A countable linearly ordered subset of $\R$ with no infinite
strict descent is well-ordered: from a nonempty subset without a least
element one could choose a strict descent recursively.

For a fixed $s$ and fixed $m$, infinitely many distinct representing
tuples would again yield a comparable pair. At least one coordinate
would be strictly larger, which would make the sums unequal. There are
only finitely many possible lengths. Finally, a decomposition into
$m$ elements of $S$ can be refined into a generator representation,
with marked cuts between its pieces. The generator representations and
the possible cuts are finite; zero pieces cause only finitely many
additional placements when $m$ is prescribed.
\end{proof}

\begin{remark}
This is the concrete positive-monoid argument required for the countably
generated support. The repository's Neumann bridge supplies two-set
product and finite-factorization statements, not by itself this entire
countable monoid construction; see Appendix~\ref{app:lean}.
\end{remark}

\subsection{Divisor coordinates and lcm multiplication}

Let $\A=\C^{\N}$, with pointwise addition and multiplication. Every
$A\in\A$ has unique divisor coordinates
\begin{equation}
 A(n)=\sum_{l\mid n}c_l,\qquad
 c_l=\sum_{d\mid l}\mu(l/d)A(d).
 \label{eq:zetamobius}
\end{equation}
These expressions are finite for each index. Equivalently,
$A=\sum_l c_l e_l$, where this notation means divisor-finite evaluation,
not a claim of uniform convergence in $n$.

\begin{proposition}[Classical lcm product]\label{prop:lcm}
If $A=\sum_i c_i e_i$ and $B=\sum_j b_j e_j$, then
\begin{equation}
 AB=\sum_l(c\star b)_l e_l,
 \qquad (c\star b)_l=\sum_{\lcm(i,j)=l}c_i b_j.
 \label{eq:lcm}
\end{equation}
The inner sum is finite. The divisor-zeta map \eqref{eq:zetamobius}
is a ring isomorphism between this lcm-coordinate ring and $\A$.
\end{proposition}
\begin{proof}
The identity $e_i(n)e_j(n)=e_{\lcm(i,j)}(n)$ follows directly from
divisibility. If $\lcm(i,j)=l$, then both $i$ and $j$ divide $l$,
so there are only finitely many such pairs. Evaluating the right side
at $n$ gives the product of the two finite divisor sums. Invertibility
of the coordinate map is M\"obius inversion, proved by
$\sum_{d\mid m}\mu(d)=\ind_{m=1}$.
\end{proof}

The construction in Proposition~\ref{prop:lcm} is classical
\cite{toth}; its role here is to preserve the arithmetic information at
action collisions.

\subsection{The formal completion}

\begin{definition}
The arithmetic Hahn ring on $S$ is
\begin{equation}
 \HH=\left\{F=\sum_{s\in S}A_s t^s:A_s\in\A\right\}.
 \label{eq:Hahn}
\end{equation}
In divisor coordinates we write
$F=\sum_{s,l}c_{s,l}e_l t^s$. Multiplication adds actions and takes lcms
of labels:
\begin{equation}
 (FG)_{u,l}=
 \sum_{\substack{s+t=u\\\lcm(i,j)=l}}c_{s,i}b_{t,j}.
 \label{eq:jointproduct}
\end{equation}
\end{definition}

Lemma~\ref{lem:support} and the finite number of divisors of $l$ make
\eqref{eq:jointproduct} a finite sum. Associativity follows by grouping
finite triple decompositions; the identity is $e_1t^0$.
For any $F$ with zero action-zero coefficient, every power $F^m$ has
actions at least $ma/2$. Hence a scalar power series $\sum_m a_mF^m$
is formally meaningful, since every exact action receives only finitely
many contributions. In particular $\exp F$, $\log(1+F)$, and
$(1+F)^{-1}$ are defined formally.

The universal arithmetic input is
\begin{equation}
 E=\sum_{d\ge2}c_de_dt^{\lambda_d}.
 \label{eq:universalE}
\end{equation}
It is legitimate for every coefficient sequence $c_d$ at the formal
level. Formal legitimacy alone does not give an analytic sum uniform in
the label. The next section supplies the missing summation condition.

\section{An action--label Banach completion}\label{sec:banach}

For a positive integer $N$ and a real $r$ with $0\le r<N$, put
\begin{equation}
 h_N(l)=\max(N,l),\qquad
 w_{N,r}(s,l)=\ee^{-s(h_N(l)-r)},
 \qquad
 \norm{F}_{N,r}=\sum_{s\in S}\sum_{l\ge1}|c_{s,l}|w_{N,r}(s,l).
 \label{eq:norm}
\end{equation}
Let $\B_{N,r}$ be the coefficient arrays of finite norm. The label
factor is essential: a divisor $l$ can occur on the diagonal only at
integers at least $l$, and its exponential decay should be measured at
that scale.

\begin{theorem}[Arithmetic Banach algebra and sheet realization]
\label{thm:banach}
The space $\B_{N,r}$, with multiplication \eqref{eq:jointproduct}, is a
commutative unital Banach algebra, and
\begin{equation}
 \norm{FG}_{N,r}\le\norm{F}_{N,r}\norm{G}_{N,r}.
 \label{eq:submult}
\end{equation}
For each $n\in\N$, the evaluation
\begin{equation}
 \mathcal E_nF(x)=\sum_{s\in S}\sum_{l\mid n}c_{s,l}\ee^{-sx}
 \label{eq:eval}
\end{equation}
converges absolutely and is bounded by $\norm{F}_{N,r}$ on
$\Re x\ge\max(N,n)-r$. It is holomorphic in the interior, and
$\mathcal E_n(FG)=\mathcal E_n(F)\mathcal E_n(G)$.
The family of all sheet evaluations, with all $n\ge1$, is faithful.
\end{theorem}
\begin{proof}
Weighted $\ell^1$ is complete because multiplication of each coordinate
by its positive weight is an isometric linear bijection to ordinary
$\ell^1(S\times\N)$. If $l=\lcm(i,j)$, then
$h_N(l)\ge h_N(i)$ and $h_N(l)\ge h_N(j)$. Since $s,t\ge0$,
\[
 w_{N,r}(s+t,l)\le w_{N,r}(s,i)w_{N,r}(t,j).
\]
Apply the triangle inequality to \eqref{eq:jointproduct}, then sum
nonnegative terms in either order. This proves \eqref{eq:submult}.
The identity has norm one.

If $l\mid n$, then $h_N(l)\le\max(N,n)$. Throughout the stated
half-plane, $|\ee^{-sx}|\le w_{N,r}(s,l)$. The norm therefore
majorizes the evaluation term by term. The same majorant proves
uniform absolute convergence, including on compact subsets of the open
half-plane, so the sum is holomorphic. Absolute convergence permits
regrouping the product by actions and lcms, giving multiplicativity.

For faithfulness, suppose every sheet function is identically zero.
Fix $n$ and let $A_s(n)=\sum_{l\mid n}c_{s,l}$. If any such coefficient
is nonzero, choose its least action $s_0$, using Lemma~\ref{lem:support}.
For large positive $x$,
\[
 \ee^{s_0x}\mathcal E_nF(x)
 =A_{s_0}(n)+\sum_{s>s_0}A_s(n)\ee^{-(s-s_0)x}.
\]
The latter sum tends to zero by dominated convergence, using absolute
convergence at one fixed sufficiently large $x_0$. Thus the identically
zero left side would tend to the nonzero number $A_{s_0}(n)$, a
contradiction. All $A_s(n)$ vanish. Apply M\"obius inversion for each
$s$ and all $n$ to obtain $c_{s,l}=0$.
\end{proof}

\begin{remark}[Faithfulness is not a tail-sequence assertion]
The complete family of analytic sheets is used in the theorem. The
nonzero arithmetic coefficient array $c_l=\mu(l)$ at action zero has
evaluation $\ind_{n=1}$. It vanishes on every diagonal integer $n\ge2$.
Thus one cannot infer formal equality merely from equality of eventual
integer sequences in this unrestricted coefficient ring.
\end{remark}

\subsection{An exact admission criterion}

\begin{theorem}[Elementary-block admission]\label{thm:admission}
For the block \eqref{eq:universalE}, the following are equivalent:
\begin{equation}
 E\in\B_{N,r}\text{ for one admissible pair }(N,r),
 \qquad
 \sum_{d\ge2}|c_d|\ee^{-ad}<\infty.
 \label{eq:admission}
\end{equation}
If these hold, then $E\in\B_{N,r}$ for every admissible pair, and for
each fixed $r\ge0$,
\begin{equation}
 \norm{E}_{N,r}\longrightarrow0\qquad(N\longrightarrow\infty).
 \label{eq:smallnorm}
\end{equation}
In particular, every polynomial-growth sequence $c_d$ is admitted.
\end{theorem}
\begin{proof}
The part with $d\le N$ is finite. For $d>N$ the norm summand is
\begin{equation}
 |c_d|\ee^{-\lambda_d(d-r)}
 =|c_d|\ee^{-ad}\ee^{a+ar-ar/d}.
 \label{eq:tailweight}
\end{equation}
The last factor is bounded above and below by positive constants
independent of $d$. This proves the equivalence and its independence of
$N,r$. For fixed $r$, each fixed-$d$ summand tends to zero as $N$ grows,
and it is bounded by $\ee^{ar}|c_d|\ee^{-a(d-1)}$. The latter sequence
is summable under \eqref{eq:admission}; dominated convergence proves
\eqref{eq:smallnorm}. Polynomial growth is dominated by any positive
exponential decay.
\end{proof}

\begin{remark}
This is an exact criterion for membership in the specified absolute
completion, not a necessary condition for the existence of each
individual $E_n$: every such sheet is a finite sum regardless of the
growth of $c_d$. The additional content is uniform control over the
expanding sheet family and its nonlinear operations.
\end{remark}

For $|c_d|\le Cd$ one has a convenient explicit majorant. Writing
$z=\ee^{-a}$, split the norm at $d=N$ and use
$\lambda_d\ge a/2$ below the split and \eqref{eq:tailweight} above it:
\begin{equation}
 \norm{E}_{N,r}\le
 \frac{CN(N+1)}2\ee^{-a(N-r)/2}
 +C\ee^{a(1+r)}
 \frac{z^{N+1}\bigl((N+1)-Nz\bigr)}{(1-z)^2}.
 \label{eq:explicitnorm}
\end{equation}
Both necklace inputs satisfy this estimate with $C=1$.

\subsection{Functional calculus, derivatives, and translation}

\begin{corollary}[Summable scalar operations]\label{cor:funcalc}
If $F\in\B_{N,r}$ and $\norm{F}_{N,r}=B<1$, then
$(1+F)^{-1}$ and $\log(1+F)$ converge in the Banach algebra. Their
truncations satisfy
\begin{align}
 \norm{(1+F)^{-1}-\sum_{m=0}^M(-F)^m}_{N,r}
 &\le\frac{B^{M+1}}{1-B},\label{eq:invbound}\\
 \norm{\log(1+F)-\sum_{m=1}^M\frac{(-1)^{m+1}}mF^m}_{N,r}
 &\le\frac{B^{M+1}}{(M+1)(1-B)}.\label{eq:logbound}
\end{align}
More generally, a scalar analytic power series converges at $F$ whenever
its radius of convergence exceeds $B$. Evaluation commutes with these
operations.
\end{corollary}
\begin{proof}
Submultiplicativity bounds the norm of each $F^m$ by $B^m$. Sum the
scalar absolute majorants. The evaluation maps have norm at most one
and are multiplicative, so they commute with the norm limits.
\end{proof}

Define the analytic-coordinate derivation by
\begin{equation}
 D(e_lt^s)=-s e_lt^s.
 \label{eq:D}
\end{equation}
It is not generally bounded on a single $\B_{N,r}$, but is bounded
with a loss of margin.

\begin{proposition}[Controlled loss of analytic margin]\label{prop:loss}
For $0\le r'<r<N$ and $k\ge0$,
\begin{equation}
 \norm{D^kF}_{N,r'}\le\frac{k!}{(r-r')^k}\norm{F}_{N,r}.
 \label{eq:loss}
\end{equation}
Translation in the analytic coordinate, defined by
$T_u(e_lt^s)=\ee^{-su}e_lt^s$, satisfies
$\norm{T_uF}_{N,r'}\le\norm{F}_{N,r}$ whenever
$|u|\le r-r'$. The label remains unchanged.
\end{proposition}
\begin{proof}
The weight ratio is $w_{N,r'}/w_{N,r}=\ee^{-s(r-r')}$.
Use $s^k\ee^{-s(r-r')}\le k!/(r-r')^k$, which follows from the
$k$th term of the exponential series. The translation assertion uses
$|\ee^{-su}|\le\ee^{s|u|}$. The Leibniz law for $D$ follows from
addition of the action indices in \eqref{eq:jointproduct}.
\end{proof}

Thus the completion is a scale of analytic algebras rather than one
Banach differential algebra with a bounded derivative. In particular,
we have not silently identified differentiation in $x$ with a shift or
variation of the arithmetic label.

\section{All-action inversion on arithmetic sheets}\label{sec:inverse}

\subsection{The exact core and the reversion equation}

For large $y>0$, the large positive solution of $\ee^{aX}/X=y$ is
\begin{equation}
 X=-\frac1a W_{-1}\!\left(-\frac ay\right),\qquad
 aX-\log X=\log y.
 \label{eq:core}
\end{equation}
Here $W_{-1}$ is the real branch of Lambert's function below $-1$.
Only the equation and the branch $X>1/a$ are needed in the proofs;
\eqref{eq:core} is an exact convenient notation, not an asymptotic
truncation. The same dominant core is already used in the repository
\cite{repo}.

Fix a label $n$ and seek $x=X+\delta$ with $f_n(x)=\ee^{aX}/X$.
After dividing the two core factors, this equation becomes
\begin{equation}
 1+E_n(X+\delta)=(1+\delta/X)\ee^{-a\delta}.
 \label{eq:root}
\end{equation}
Set
\begin{equation}
 J_X(u)=\frac{1-(1+u/X)\ee^{-au}}u,
 \qquad p=J_X(0)=a-\frac1X.
 \label{eq:J}
\end{equation}
The removable value at $u=0$ is understood. Introducing a coupling
parameter $z$ gives the exact implicit equation
\begin{equation}
 \delta J_X(\delta)+zE_n(X+\delta)=0.
 \label{eq:coupling}
\end{equation}
We expand in $z$ at zero and then set $z=1$. The coupling order counts
how many primitive sectors participate; it is not an action cutoff.

\subsection{The algebraic reversion identity}

\begin{lemma}[Lagrange coefficient identity]\label{lem:lagrange}
For $\delta=zF(\delta)$ with $\delta(0)=0$, over any commutative
$\Q$-algebra, the unique formal solution satisfies
\begin{equation}
 [z^m]G(\delta)=\frac1m[u^{m-1}]G'(u)F(u)^m,
 \qquad m\ge1.
 \label{eq:lagrange}
\end{equation}
The coefficients exist without an invertibility assumption on $F(0)$.
\end{lemma}
\begin{proof}
The coefficient of $z^m$ in $\delta$ is determined recursively by lower
coefficients, proving existence and uniqueness. When $F(0)$ is an
indeterminate unit, formal residue substitution $z=u/F(u)$ gives
\[
 [z^m]G(\delta)
 =\Res_{u=0}G(u)\left(\frac{F(u)^m}{u^{m+1}}
       -\frac{F(u)^{m-1}F'(u)}{u^m}\right)du.
\]
The residue of the derivative of $G(u)F(u)^m/u^m$ is zero.
Rearranging this identity gives \eqref{eq:lagrange}. For a fixed $m$,
both sides are polynomial expressions with rational coefficients in
finitely many coefficients of $F$ and $G$. An identity proved after
localizing the universal polynomial ring at its nonzero indeterminate
$F(0)$ is therefore already an identity in that polynomial ring, and
specializes to any commutative $\Q$-algebra. This includes rings with
idempotents and zero divisors.
\end{proof}

Apply the lemma to $F(u)=-E_n(X+u)/J_X(u)$ and $G(u)=u$. Then
\begin{align}
 \delta_n(X)&=\sum_{m\ge1}\delta_{m,n}(X),\label{eq:deltasum}\\
 \delta_{m,n}(X)
 &=\frac{(-1)^m}{m}[u^{m-1}]
       \left(\frac{E_n(X+u)}{J_X(u)}\right)^m.
 \label{eq:deltam}
\end{align}
At first this is a formal expression. Expanding the power of $E_n$
reveals the arithmetic action structure:
\begin{align}
 \delta_n(X)
 &=\sum_{m\ge1}\ \sum_{d_1,\ldots,d_m\ge2}
 C_m(s,X)\left(\prod_{i=1}^m c_{d_i}\right)
 e_{\lcm(d_1,\ldots,d_m)}(n)\ee^{-sX},\label{eq:allactions}\\
 s&=\lambda_{d_1}+\cdots+\lambda_{d_m},\nonumber\\
 C_m(s,X)&=\frac{(-1)^m}{m}[u^{m-1}]\ee^{-su}J_X(u)^{-m}.
 \label{eq:Cm}
\end{align}
The inner sum is over ordered tuples. For each exact $s$ it is finite,
including the possible $m$, by Lemma~\ref{lem:support}. The coefficients
$C_m(s,X)$ are rational functions of $X$ and $a$, with $s/a\in\Q$,
whenever $p\ne0$. Thus \eqref{eq:allactions} is a formal arithmetic
Hahn series with these rational functions as an additional coefficient
field. No fractional divisibility predicate has appeared.

\subsection{A quantitative convergence theorem}

\begin{theorem}[Absolute inverse transport with an explicit tail]
\label{thm:inverse}
Let $a>0$, $X\ge4/a$, and $0<\rho\le1/(8a)$. Fix $n$ and any
complex coefficients $c_d$ on its divisors. Define $A_n(X)$ by
\eqref{eq:An} without using its final upper bound, and set
\begin{equation}
 \theta=\frac{2\ee^{a\rho}A_n(X)}{a\rho}.
 \label{eq:theta}
\end{equation}
If $\theta<1$, equation \eqref{eq:root} has a unique solution
$\delta$ in $|\delta|<\rho$. It is the sum \eqref{eq:deltasum},
and both that sum and its full ordered-tuple expansion
\eqref{eq:allactions} converge absolutely. For every $M\ge0$,
\begin{equation}
 \abs{\delta-\sum_{m=1}^M\delta_{m,n}(X)}
 \le\frac{\rho\theta^{M+1}}{(M+1)(1-\theta)}.
 \label{eq:inverseerror}
\end{equation}
The same expression bounds the sum of the absolute values of all
omitted tuple contributions.
\end{theorem}
\begin{proof}
First bound the denominator on the closed disk. Integrating the
derivative of $(1+u/X)\ee^{-au}$ gives
\begin{equation}
 J_X(u)=\int_0^1\ee^{-atu}\left(p+\frac{atu}{X}\right)dt.
 \label{eq:Jintegral}
\end{equation}
Here $3a/4\le p<a$ and $\rho/X\le1/32$. Consequently,
\[
 |J_X(u)-p|
 \le a(\ee^{1/8}-1)+\frac a{32}\ee^{1/8}
 \le\frac a7+\frac a{28}=\frac{5a}{28}.
\]
The last inequality follows from $\ee^{1/8}\le(1-1/8)^{-1}=8/7$.
It follows that $|J_X(u)|\ge4a/7>a/2$ throughout the disk. In
particular, $J_X$ has no zero there.

On $|u|=\rho$,
\[
 |E_n(X+u)|\le\ee^{a\rho}A_n(X),\qquad
 |uJ_X(u)|>a\rho/2.
\]
For $|z|<1/\theta$, Rouch\'e's theorem applied to
$uJ_X(u)+zE_n(X+u)$ gives exactly one zero in the disk, counted with
multiplicity. It is therefore a simple zero. The local holomorphic
implicit solutions patch, by uniqueness, into one holomorphic function
$\delta(z)$ on $|z|<1/\theta$ with $\delta(0)=0$.
At $z=1$ it solves \eqref{eq:root}. Its Taylor coefficients are
\eqref{eq:deltam} by Lemma~\ref{lem:lagrange}.

For absolute summation at the finer tuple level, the Cauchy coefficient
bound in \eqref{eq:Cm} is
\begin{equation}
 |C_m(s,X)|\le\frac1{m\rho^{m-1}}
 \left(\frac2a\right)^m\ee^{s\rho}.
 \label{eq:Cbound}
\end{equation}
For a fixed label only divisor tuples survive. Summing the absolute
values of their contributions and using $\lambda_d<a$ gives
\[
 \sum_{\substack{d_1,\ldots,d_m\mid n\\d_i\ge2}}
 |C_m(s,X)|\prod_i|c_{d_i}|\ee^{-sX}
 \le\frac\rho m\left(
 \frac{2\ee^{a\rho}A_n(X)}{a\rho}\right)^m
 =\frac\rho m\theta^m.
\]
Summing this bound from $m=M+1$ to infinity proves
\eqref{eq:inverseerror}, since $1/m\le1/(M+1)$ on that range.
It also proves absolute convergence before grouping actions, so all
regroupings used in \eqref{eq:allactions} are legitimate.
\end{proof}

\begin{remark}[The vanishing-input case]
If $A_n(X)=0$, then $E_n$ is zero, the unique root in the disk is
$\delta=0$, and all coefficients and error bounds vanish. The proof
using $1/\theta$ is read separately in this case.
\end{remark}

\begin{corollary}[Uniformity supplied by the Banach completion]
\label{cor:uniforminverse}
Suppose $\sum_{d\ge2}|c_d|\ee^{-ad}<\infty$, and fix
$0<\rho\le1/(8a)$. For every sufficiently large integer $N$, inversion
is uniformly convergent for all labels $n$ and real core coordinates
\begin{equation}
 X\ge\max(N,n)-\rho,\qquad X\ge4/a.
 \label{eq:uniformdomain}
\end{equation}
An admissible common majorant in \eqref{eq:inverseerror} is
\begin{equation}
 \Theta_N=\frac{2}{a\rho}\norm{E}_{N,2\rho}<1.
 \label{eq:ThetaN}
\end{equation}
\end{corollary}
\begin{proof}
Choose $N>2\rho$. For $|u|\le\rho$ and a divisor $d\mid n$,
\[
 |c_d\ee^{-\lambda_d(X+u)}|
 \le |c_d|\ee^{-\lambda_d(h_N(d)-2\rho)}.
\]
The total absolute majorant on the circle is at most
$\norm{E}_{N,2\rho}$. Repeat the proof of Theorem~\ref{thm:inverse}
with this majorant instead of $\ee^{a\rho}A_n(X)$, including the
ordered-tuple estimate. The norm tends to zero by
Theorem~\ref{thm:admission}; hence $\Theta_N<1$ eventually.
\end{proof}

This corollary makes the connection between the global arithmetic
completion and the local inverse theorem explicit. The inverse's
coefficients depend rationally on $X$; the assertion is a uniformly
convergent family of such inverses, not an unannounced claim that
variable-$X$ coefficients belong to the original constant-coefficient
Banach space.

\subsection{An exact inverse on the range}

\begin{theorem}[Arithmetic range inverse]\label{thm:range}
Let $c_d$ be real and satisfy the admission condition
\eqref{eq:admission}. Put
\begin{equation}
 y_n=\frac{\ee^{an}}n\bigl(1+E_n(n)\bigr).
 \label{eq:yn}
\end{equation}
For all sufficiently large $n$, this is positive, and its large core
inverse $X_n>1/a$ satisfies $X_n-n\to0$. The convergent series
\eqref{eq:allactions}, evaluated with label $n$ and core $X_n$, obeys
\begin{equation}
 n=X_n+\sum_{m\ge1}\delta_{m,n}(X_n).
 \label{eq:rangeinverse}
\end{equation}
For necklaces and primitive necklaces,
\begin{equation}
 X_n-n=\bigO_a\!\left(n\ee^{-an/2}\right),
 \label{eq:Xn}
\end{equation}
and the explicit bound \eqref{eq:inverseerror} applies with
$\theta\le 2\ee^{a\rho}n\ee^{-aX_n/2}/(a\rho)$.
\end{theorem}
\begin{proof}
By Theorem~\ref{thm:admission} and diagonal evaluation,
$|E_n(n)|\le\norm{E}_{n,0}\to0$. In particular $1+E_n(n)>0$
for large $n$. Write $H(x)=ax-\log x$. Its derivative is
$a-1/x\ge a/2$ for $x\ge2/a$, and
\[
 H(X_n)-H(n)=\log(1+E_n(n)).
\]
For large $n$ the large core solution is also at least $2/a$, and the
mean value theorem gives
$|X_n-n|\le(2/a)|\log(1+E_n(n))|\to0$.
For the two counting inputs, \eqref{eq:An} gives
\eqref{eq:Xn} after using $|\log(1+t)|\le2|t|$ for $|t|\le1/2$.

Fix a permitted $\rho$. Eventually $|X_n-n|<\rho$ and the uniform
inverse theorem applies. The number $\delta=n-X_n$ already satisfies
\eqref{eq:root}, because $f_n(n)=y_n=\ee^{aX_n}/X_n$. Its size is
less than $\rho$, so uniqueness identifies it with the convergent
Lagrange series. The stated explicit bound follows from
\eqref{eq:An}.
\end{proof}

\begin{remark}[What this does and does not invert]
Equation \eqref{eq:rangeinverse} is an identity on the exact sequence
range and a convergent local inverse on each specified sheet. It does
not require choosing a global interpolation. It also does not by itself
solve the problem of determining an unknown label from an arbitrary
real target: the label occurs in the sheet data. On the exact range,
\eqref{eq:Xn} already shows that nearest-integer rounding of $X_n$
recovers $n$ eventually. A practical threshold problem can use candidate
labels and exact comparisons, as in the repository; no new general
threshold theorem is claimed here.
\end{remark}

\section{Coefficients, collisions, and the first post-cut sector}
\label{sec:firstcut}

\subsection{The first three homogeneous kernels}

Expand $J_X(u)=p+j_1u+j_2u^2+\cdots$. Directly from
\eqref{eq:J},
\begin{equation}
 j_1=-\frac{a^2}2+\frac aX,
 \qquad j_2=\frac{a^3}6-\frac{a^2}{2X}.
 \label{eq:js}
\end{equation}
The first kernels from \eqref{eq:Cm} are
\begin{align}
 C_1(s,X)&=-\frac1p,\label{eq:Cfirst}\\
 C_2(s,X)&=-\frac{s}{2p^2}-\frac{j_1}{p^3},\label{eq:Csecond}\\
 C_3(s,X)&=-\frac{s^2}{6p^3}-\frac{s j_1}{p^4}
          -\frac{2j_1^2}{p^5}+\frac{j_2}{p^4}.
 \label{eq:Cthird}
\end{align}
For example, the second expression comes from the coefficient of $u$
in $\ee^{-su}J_X(u)^{-2}/2$. The third follows from the coefficient
of $u^2$ in the same expansion with exponent $-3$, multiplied by
$-1/3$. Equivalently,
\begin{equation}
 \delta_{1,n}=-\frac{E_n(X)}p,
 \qquad
 \delta_{2,n}=\frac{E_n(X)E'_n(X)}{p^2}
              -\frac{j_1E_n(X)^2}{p^3}.
 \label{eq:deltatwo}
\end{equation}
This displays why differentiating the small sectors while holding their
arithmetic indicators fixed is necessary.

\subsection{The arithmetic content of action collisions}

At action $a$, the only representation is
$\lambda_2+\lambda_2=a$. At action $4a/3$, the unordered two-term
representations are $(2,6)$ and $(3,3)$. At action $3a/2$, they are
$(3,6)$ and $(4,4)$, and there is also the three-term representation
$(2,2,2)$. There are no one-term representations at these actions.
The completeness of these short lists can be checked by solving the
corresponding two-unit-fraction equations, or by the finite algorithm in
Section~\ref{sec:algorithm}.

For $c_d=\varphi(d)$, the coefficients multiplying the indicated
exponential monomial are
\begin{align}
 [\ee^{-aX}]\delta_n&=C_2(a,X)e_2(n),\nonumber\\
 [\ee^{-4aX/3}]\delta_n
 &=C_2(4a/3,X)\bigl(4e_3(n)+4e_6(n)\bigr),\label{eq:phicollisions}\\
 [\ee^{-3aX/2}]\delta_n
 &=C_2(3a/2,X)\bigl(8e_6(n)+4e_4(n)\bigr)
   +C_3(3a/2,X)e_2(n).\nonumber
\end{align}
The bracket notation here means an exact action coefficient in the
formal series, not extraction from a one-variable ordinary power series.
For $c_d=\mu(d)$ the corresponding expressions are
\begin{align}
 [\ee^{-aX}]\delta_n&=C_2(a,X)e_2(n),\nonumber\\
 [\ee^{-4aX/3}]\delta_n
 &=C_2(4a/3,X)\bigl(e_3(n)-2e_6(n)\bigr),\label{eq:mucollisions}\\
 [\ee^{-3aX/2}]\delta_n
 &=-2C_2(3a/2,X)e_6(n)-C_3(3a/2,X)e_2(n).\nonumber
\end{align}
For instance the two ordered tuples $(2,6)$ and $(6,2)$ contribute
$2c_2c_6e_6$, whereas $(3,3)$ contributes $c_3^2e_3$.
Keeping only the total action would lose this distinction.

\subsection{Cancellation at the accumulation boundary}

Substitute $s=a$ into \eqref{eq:Csecond}. A cancellation gives the
exact identity
\begin{equation}
 C_2(a,X)=-\frac{a}{2X(a-1/X)^3}.
 \label{eq:boundarykernel}
\end{equation}
Its limit as $X\to\infty$ is zero. Thus an action at the first
nonlinear boundary is present, but its amplitude is one algebraic order
smaller than a naive nonresonant estimate suggests.

\begin{theorem}[First nonlinear sector beyond the primitive cut]
\label{thm:firstcut}
Fix $a=\log q>0$ and either $c_d=\varphi(d)$ or $c_d=\mu(d)$.
Fix $0<\rho\le1/(8a)$. Uniformly for real $X$ with $|X-n|\le\rho$
as $n\to\infty$, the inverse displacement of
Theorem~\ref{thm:inverse} satisfies
\begin{equation}
 \delta_n(X)+\frac{E_n(X)}{a-1/X}
 =-\frac{a\,e_2(n)}{2X(a-1/X)^3}\ee^{-aX}
 +\bigO_a\!\left(n^2\ee^{-7aX/6}\right).
 \label{eq:firstcut}
\end{equation}
In particular, on even labels,
\begin{equation}
 \delta_n(X)+\frac{E_n(X)}{a-1/X}
 \sim-\frac{a}{2X(a-1/X)^3}\ee^{-aX}.
 \label{eq:firstequiv}
\end{equation}
The same statements hold at the exact range core $X=X_n$.
\end{theorem}
\begin{proof}
The first homogeneous term is exactly $-E_n(X)/p$, so it is removed
in the left side. The unique quadratic tuple with action $a$ is
$(2,2)$. In both counting cases $c_2^2=1$, and its label is $e_2$.
Its coefficient is \eqref{eq:boundarykernel}.

Every other quadratic tuple has at least one index at least three,
so its action is at least
$\lambda_2+\lambda_3=7a/6$. The kernel $C_2(s,X)$ is uniformly
bounded for $s\in[a,2a]$ and large $X$. The sum of absolute products
of the divisor coefficients is at most $n^2$, by
\eqref{eq:arithbounds}. The remaining quadratic terms are therefore
$\bigO_a(n^2\ee^{-7aX/6})$.

For the remaining estimates, use the fixed disk radius $1/(8a)$ in
Theorem~\ref{thm:inverse}; uniqueness identifies its root with the root
in any smaller permitted disk for large $n$. This makes the constants
depend only on $a$. For $m\ge3$, its absolute tail is
$\bigO_a(A_n(X)^3)=\bigO_a(n^3\ee^{-3aX/2})$, uniformly in the
stated range, since $\theta\to0$. This is absorbed by the preceding
bound: their ratio is $n\ee^{-aX/3}\to0$. This proves
\eqref{eq:firstcut}. On even labels the ratio of the error to the
absolute leading term is
$\bigO_a(Xn^2\ee^{-aX/6})\to0$, giving
\eqref{eq:firstequiv}. The range assertion follows from
\eqref{eq:Xn}.
\end{proof}

\begin{remark}[Why the full primitive block must be retained]
The subtracted term $E_n(X)/p$ includes every divisor of $n$. For
necklaces, if a divisor cutoff omits $d=n$, it omits the positive term
$\varphi(n)\ee^{-aX+aX/n}$. When $X=n+\bigO(1)$, this term is at
least a positive constant times $\ee^{-aX}$, already larger than the
$X^{-1}\ee^{-aX}$ effect in \eqref{eq:firstcut}. The theorem cannot
be recovered by subtracting a fixed finite primitive prefix and then
reading off the quadratic sector. This is exactly where summation
across the accumulating block becomes mathematically necessary.
\end{remark}

\section{An all-order resonance law}\label{sec:resonance}

The cancellation in \eqref{eq:boundarykernel} is not isolated. We now
transform the coefficient kernel so that its entire pure-exponential
limit, and the first correction caused by the factor $x^{-1}$, are
explicit.

For complex $\alpha$ and $m\ge0$, let
$\binom{\alpha}{m}=\alpha(\alpha-1)\cdots(\alpha-m+1)/m!$.
Define the analytic germ
\begin{equation}
 B(w)=\frac{(1+w)\log(1+w)}w,
 \qquad B(0)=1,
 \qquad \eps=\frac1{aX}.
 \label{eq:Bkernel}
\end{equation}
The logarithm is the germ at $w=0$.

\begin{theorem}[Exact transformed kernel and curvature lifting]
\label{thm:resonance}
For every $m\ge1$, $s$, and $X$ with $p\ne0$, put $\alpha=s/a$.
The kernel \eqref{eq:Cm} has the exact coefficient representation
\begin{equation}
 C_m(s,X)=-\frac1{am}[w^{m-1}]
  (1+w)^{\alpha-1}\bigl(1-\eps B(w)\bigr)^{-m}.
 \label{eq:transformedC}
\end{equation}
For fixed $m,\alpha$, as $X\to\infty$,
\begin{equation}
 C_m(s,X)=-\frac1{am}\binom{\alpha-1}{m-1}
 -\frac1{a^2X}\frac{d}{d\alpha}\binom{\alpha}{m}
 +\bigO_{a,m,\alpha}(X^{-2}).
 \label{eq:Cexpansion}
\end{equation}
At every integer $k$ with $1\le k<m$, the first term vanishes and
\begin{equation}
 C_m(ka,X)=
 \frac{(-1)^{m-k}}{a^2X\,m\binom{m-1}{k}}
 +\bigO_{a,m,k}(X^{-2}).
 \label{eq:resonant}
\end{equation}
Whenever the action $ka$ has an $m$-tuple representation in the
arithmetic support, this is its universal homogeneous kernel. The
arithmetic sum of tuple amplitudes can, independently, vanish.
\end{theorem}
\begin{proof}
It is useful to derive the identity for a finite collection of sector
amplitudes $b_d=c_de_d(n)\ee^{-\lambda_dX}$, treated as independent
commuting quantities. In \eqref{eq:root} with coupling $z$, put
$v=\ee^{-a\delta}$ and $w=v-1$. Since
$\delta=-a^{-1}\log v$, the equation becomes
\[
 v\bigl(1-\eps\log v\bigr)
 =1+z\sum_d b_d v^{\lambda_d/a}.
\]
Equivalently,
\begin{equation}
 w\bigl(1-\eps B(w)\bigr)
 =z\sum_d b_d(1+w)^{\lambda_d/a}.
 \label{eq:wequation}
\end{equation}
Use Lemma~\ref{lem:lagrange} on this equation with
$G(w)=-a^{-1}\log(1+w)$. The contribution of an ordered tuple of
length $m$ is precisely the right side of
\eqref{eq:transformedC}, since $G'(w)=-1/(a(1+w))$ and the tuple's
exponents sum to $\alpha$. The identity is a universal coefficient
identity, so it also proves equality with \eqref{eq:Cm}. Only finitely
many Taylor coefficients are involved at fixed $m$.

Expand the last factor in \eqref{eq:transformedC} at $\eps=0$:
\[
 (1-\eps B(w))^{-m}=1+m\eps B(w)+\bigO(\eps^2).
\]
The constant term in $\eps$ gives the first term of
\eqref{eq:Cexpansion}. For the linear term, the definition of $B$
gives
\begin{align*}
 [w^{m-1}](1+w)^{\alpha-1}B(w)
 &=[w^m](1+w)^\alpha\log(1+w)\\
 &=\frac{d}{d\alpha}\binom{\alpha}{m}.
\end{align*}
Since $\eps/a=1/(a^2X)$, this proves \eqref{eq:Cexpansion}.
The remainder estimate is valid for fixed $m,\alpha$: the relevant
finite coefficient extraction is analytic in $\eps$ near zero.

At $\alpha=k\in\{1,\ldots,m-1\}$,
$\binom{k-1}{m-1}=0$. Differentiating the polynomial
$\binom{\alpha}{m}$ at its simple root $k$ gives
\[
 \left.\frac{d}{d\alpha}\binom{\alpha}{m}\right|_{\alpha=k}
 =\frac{k!(-1)^{m-1-k}(m-1-k)!}{m!}
 =\frac{(-1)^{m-1-k}}{m\binom{m-1}{k}}.
\]
Substitution proves \eqref{eq:resonant}.
\end{proof}

For the necklace action set, an $m$-tuple has
$m/2\le\alpha<m$. Thus only integers in this interval can occur as
actual resonant actions, even though the kernel identity holds for all
$k$ in the theorem. For $m=2$, $k=1$, formula \eqref{eq:resonant}
reads $-1/(2a^2X)+\bigO(X^{-2})$, in agreement with the exact
expression \eqref{eq:boundarykernel}.

\subsection{The pure-exponential comparison}

If the factor $x^{-1}$ is removed from the forward core, then
$J_X$ is replaced by $J_\infty(u)=(1-\ee^{-au})/u$ and the kernel is
exactly
\begin{equation}
 C_m^{\mathrm{exp}}(s)=-\frac1{am}\binom{s/a-1}{m-1}.
 \label{eq:purekernel}
\end{equation}
This follows either by putting $\eps=0$ in \eqref{eq:wequation},
or by the same change of coordinate for the pure-exponential equation.
The integer zeros in \eqref{eq:purekernel} are therefore exact
cancellations of the entire homogeneous kernel. The logarithmic term
in the phase $ax-\log x$ removes these zeros at order $X^{-1}$.
This is the sense in which curvature lifts the resonance.

\subsection{General power-law factors}

Consider instead
\begin{equation}
 f_n^{(b)}(x)=\ee^{ax}x^{-b}\bigl(1+E_n(x)\bigr),
 \qquad b\in\R\text{ fixed},
 \label{eq:bcore}
\end{equation}
near a large positive $X$, with the local analytic branch of $x^{-b}$.
Replace $J_X$ by
\begin{equation}
 J_X^{(b)}(u)=\frac{1-(1+u/X)^b\ee^{-au}}u,
 \qquad J_X^{(b)}(0)=a-b/X.
 \label{eq:Jb}
\end{equation}
For fixed $b$ this is uniformly nonzero on a sufficiently small fixed
disk for all large $X$, because it converges there to $J_\infty$.
Thus the same Rouch\'e and coefficient arguments provide an inverse
series and a geometric tail, with the lower bound for $J_X^{(b)}$
substituted in place of $a/2$.

\begin{corollary}[Universal first curvature coefficient]
\label{cor:bcore}
Let $C_m^{(b)}(s,X)$ be defined by \eqref{eq:Cm} with $J_X^{(b)}$.
For fixed $a,b,m,\alpha=s/a$,
\begin{equation}
 C_m^{(b)}(s,X)=
 -\frac1{am}\binom{\alpha-1}{m-1}
 -\frac{b}{a^2X}\frac{d}{d\alpha}\binom{\alpha}{m}
 +\bigO(X^{-2}).
 \label{eq:bexpansion}
\end{equation}
At an integer resonance $1\le k<m$ this gives
\begin{equation}
 C_m^{(b)}(ka,X)=
 \frac{(-1)^{m-k}b}{a^2X\,m\binom{m-1}{k}}
 +\bigO(X^{-2}).
 \label{eq:bresonance}
\end{equation}
\end{corollary}
\begin{proof}
The change $v=1+w=\ee^{-a\delta}$ now gives
\[
 (1+w)\bigl(1-\eps\log(1+w)\bigr)^b
 =1+z\sum_d b_d(1+w)^{\lambda_d/a}.
\]
After subtracting one and dividing by $w$, the factor multiplying $w$
is
\[
 K_b(w,\eps)=
 \frac{(1+w)(1-\eps\log(1+w))^b-1}{w}
 =1-b\eps B(w)+\bigO(\eps^2).
\]
All expressions are analytic germs, including at $w=0$.
Lagrange extraction now contains $K_b^{-m}$, whose expansion is
$1+mb\eps B(w)+\bigO(\eps^2)$. Repeat the coefficient calculation
in Theorem~\ref{thm:resonance}. The simple-root derivative formula
then gives \eqref{eq:bresonance}.
\end{proof}

The bounds in \eqref{eq:Cexpansion} and \eqref{eq:bexpansion} are for
fixed homogeneous order. They are not uniform when $m$ grows with $X$.
The joint large-order problem is genuinely additional, as noted in
Section~\ref{sec:questions}.

\section{Obstructions to stronger interpretations}\label{sec:obstructions}

The limitations below are structural statements with proofs, not missing
steps concealed inside the inversion theorem.

\subsection{Indicators cannot live faithfully in a field}

\begin{proposition}[Idempotent obstruction]\label{prop:idempotent}
No unital ring embedding of the arithmetic coefficient ring $\A$ into a
field exists. Every derivation on a commutative ring containing the
indicators annihilates each $e_d$. Consequently the shift
$A(n)\mapsto A(n+1)$ cannot be the exponential of a derivation that
preserves the indicator ring.
\end{proposition}
\begin{proof}
The indicator $e_2$ is a nonzero idempotent different from one, and
$e_2(1-e_2)=0$ with both factors nonzero. A field has no such pair,
so an injective ring map is impossible.

For any idempotent $e$, the Leibniz rule gives
$D(e)=D(e^2)=2eD(e)$, whence $(1-2e)D(e)=0$.
But $(1-2e)^2=1$, so $D(e)=0$. Thus every iterate of $D$ annihilates
$e_d$, and any meaningful exponential of $D$ fixes $e_d$. The integer
shift sends $e_2$ to $1-e_2$, so it cannot be such an exponential.
\end{proof}

The arithmetic ring therefore is not an ordinary transseries field with
an enlarged set of constant symbols. A difference operator acting on
labels can be studied separately, but it is not supplied by analytic
translation in $x$. This is why the frozen-label convention is necessary
rather than cosmetic.

\subsection{A compact profinite parameter does not repair summation}

Let $\widehat{\mathbb Z}$ denote the profinite completion of the
integers. Each single indicator $e_d$ extends continuously to it, because
it is the characteristic function of the clopen set $d\widehat{\mathbb Z}$.
It is tempting to sum the primitive block at fixed $x$ and regard the
result as a continuous function of a profinite label. The following
proposition shows that this fails for both counting inputs.

\begin{proposition}[Fixed-coordinate profinite obstruction]
\label{prop:profinite}
Fix a real $x>0$. Neither of the functions
$n\mapsto E_n(x)$ for $c_d=\varphi(d)$ or $c_d=\mu(d)$ is the
restriction to positive integers of a continuous complex-valued function
on $\widehat{\mathbb Z}$.
\end{proposition}
\begin{proof}
Put $b=ax>0$. In the totient case, $\lambda_d<a$ and positivity imply
\[
 E_n(x)\ge\ee^{-b}\sum_{d\mid n,\ d\ge2}\varphi(d)
 =\ee^{-b}(n-1).
\]
Thus the restriction is unbounded, whereas every continuous function on
the compact space $\widehat{\mathbb Z}$ is bounded.

For the M\"obius case, including the $d=1$ term gives
\[
 1+E_n(x)=\ee^{-b}\sum_{d\mid n}\mu(d)\ee^{b/d}.
\]
For $n>1$, expand the exponential and interchange with the finite
divisor sum. The constant coefficient vanishes, and
\begin{equation}
 \sum_{d\mid n}\mu(d)\ee^{b/d}
 =\sum_{k\ge1}\frac{b^k}{k!}
       \prod_{p\mid n}(1-p^{-k}).
 \label{eq:muprof}
\end{equation}
Every product is between zero and one, so
$1+E_n(x)\le1-\ee^{-b}$, or $E_n(x)\le-\ee^{-b}$, for all $n>1$.
On the other hand, $E_1(x)=0$. The positive integers $n_j=1+j!$
converge to 1 in the profinite topology: every fixed modulus divides
$j!$ eventually. The values $E_{n_j}(x)$ stay separated from zero,
contradicting continuity at 1.
\end{proof}

The weighted sheet domain avoids this contradiction by requiring the
analytic coordinate to grow with the label. It is not a product of a
fixed analytic domain with all of $\widehat{\mathbb Z}$. Other weighted
or partially defined profinite constructions remain possible, but they
must respect the obstruction just proved.

\subsection{Bounded-action extraction versus homogeneous truncation}

Theorem~\ref{thm:inverse} truncates at a homogeneous order $M$. Each
retained term sums the entire primitive divisor block, and the omitted
terms carry at least $M+1$ primitive factors. This is different from
listing all actions below a fixed numerical cutoff. Before action $a$
alone there are infinitely many primitive actions in the universal
series. No theorem above converts that infinite list into a finite grid.

At a fixed label $n$ there are only finitely many primitive divisors,
so a computation is finite. Uniformity over all $n$ is supplied by the
weighted norm, not by silently interchanging a label-dependent finite
sum with a globally finite action list. Equally, summability in this
article means ordinary absolute convergence of the specified families;
no Borel summation, alien derivative, or Stokes automorphism is asserted.

\section{Finite algorithms and reproducible audits}\label{sec:algorithm}

\subsection{Exact membership and a fixed-action coefficient algorithm}

\begin{proposition}[A finite action search]\label{prop:actionalgorithm}
For a given rational $\alpha\ge0$, it is decidable by a finite exact
search whether $a\alpha\in S$, and all generator representations can
be enumerated. Hence every exact action coefficient in
\eqref{eq:allactions} is computable from a finite list of the input
coefficients $c_d$.
\end{proposition}
\begin{proof}
For $\alpha=0$ take the empty tuple. Otherwise a representation of
length $m$ must have $m\le\lfloor2\alpha\rfloor$ and obey
\begin{equation}
 \frac1{d_1}+\cdots+\frac1{d_m}=m-\alpha>0.
 \label{eq:egyptian}
\end{equation}
Search only sorted denominators. At a recursive step with $k$ terms
remaining, target sum $t>0$, and smallest permitted denominator $L$,
the next denominator lies in the finite interval
\begin{equation}
 \max\!\left(L,\left\lceil\frac1t\right\rceil\right)
 \le d\le\left\lfloor\frac{k}{t}\right\rfloor.
 \label{eq:denomrange}
\end{equation}
The lower bound follows from $1/d\le t$. The upper bound follows
because the remaining sorted denominators are at least $d$, so the
sum of the $k$ reciprocals is at most $k/d$. Recurse with target
$t-1/d$, $k-1$ terms, and lower bound $d$. Nonpositive targets are
rejected unless no terms remain and the target is zero. For $k=1$
the sole possible denominator is $1/t$, which must be an integer at
least $L$. The depth and every branching interval are finite.

Restore the ordered tuples by permutation. Equivalently, a sorted tuple
with multiplicities $r_d$ receives the factor $m!/\prod_d r_d!$.
The action coefficient follows by multiplying the corresponding
$c_d$, using the lcm of the denominators for the label, and applying
\eqref{eq:Cm}. No bound on the running time polynomial in the bit
length of $\alpha$ is claimed.
\end{proof}

\subsection{Efficient homogeneous jets for a known label}

Enumerating all tuples is unnecessary for numerical inversion on one
sheet. Given the divisors of $n$, form the weighted moments
\begin{equation}
 W_k=\sum_{d\mid n,\ d\ge2}c_d\lambda_d^k\ee^{-\lambda_dX},
 \qquad 0\le k<M.
 \label{eq:Wk}
\end{equation}
The Taylor coefficients required by \eqref{eq:deltam} are
\begin{align}
 [u^k]E_n(X+u)&=\frac{(-1)^k}{k!}W_k,\label{eq:Ejet}\\
 [u^k]J_X(u)&=-\frac{(-a)^{k+1}}{(k+1)!}
              -\frac{(-a)^k}{Xk!}.\label{eq:Jjet}
\end{align}
Divide the two jets modulo $u^M$, obtaining
$F(u)=E_n(X+u)/J_X(u)$. Successively multiply by $F$ modulo $u^M$
and read $(-1)^m[u^{m-1}]F(u)^m/m$ for $m=1,\ldots,M$.

With naive polynomial arithmetic, the moment construction and these
jet operations take $\bigO(M\tau(n)+M^3)$ scalar arithmetic
operations once the divisors and exponentials are available. This is
not a bit-complexity or integer-factorization bound. The supplied
implementation uses trial division to enumerate divisors, and can be
replaced by a faster divisor provider without changing the coefficient
calculation. Large precision is important when checking an exponentially
small residual by subtracting two nearly equal inverse values.

For a genuinely machine-certified numerical enclosure, every input
bound and polynomial operation must use outward-rounded intervals.
The analytic inequality \eqref{eq:inverseerror} is rigorous; the
recorded floating-point evaluations of it are diagnostic values, not
interval certificates.

\subsection{What the accompanying program checks}

The file \texttt{verification/verify.py} is self-contained except for
\texttt{mpmath}. Its exact portion uses Python integer arithmetic and
\texttt{fractions.Fraction}. It performs 964 successful exact assertions:
it checks finite lcm-convolution identities, divisor identities up to
$n=300$, substitution of the rational formal inverse through coupling
order seven, the integer-root derivative formula through $m=14$, four
specified action-collision lists, and 80 finite-$X$ identities comparing
the two coefficient coordinates in \eqref{eq:Cm} and
\eqref{eq:transformedC}.

The numerical portion uses 400 decimal digits. Exact integer necklace
or primitive-necklace counts provide the targets, not floating-point
forward approximations. For twelve combinations of counting type,
$q$, and $n$, it computes the large Lambert core $X$, compares the first
four homogeneous inverse sums with the exact displacement $n-X$, and
checks all 48 instances of \eqref{eq:inverseerror}. It also records
ratios testing \eqref{eq:firstequiv}. The generated CSV and JSON files
are included so that the reported checks can be rerun.

\begin{table}[htbp]
\centering
\caption{Binary-necklace range inversion. Actual absolute error of the
first $M$ homogeneous terms, and the evaluated bound from
\eqref{eq:inverseerror}. These are high-precision diagnostics, not
outward-rounded intervals. Values are generated by the supplied program.}
\label{tab:errors}
\begin{tabular}{@{}rrcc@{}}
\toprule
$n$&$M$&Absolute error&Analytic-bound value\\
\midrule
% BEGIN GENERATED VERIFICATION TABLE
20 & 1 & $1.708\times10^{-7}$ & $3.853\times10^{-5}$ \\
20 & 2 & $4.958\times10^{-11}$ & $5.255\times10^{-7}$ \\
20 & 4 & $2.195\times10^{-18}$ & $1.32\times10^{-10}$ \\
60 & 1 & $1.671\times10^{-20}$ & $2.581\times10^{-17}$ \\
60 & 2 & $4.953\times10^{-29}$ & $2.912\times10^{-25}$ \\
60 & 4 & $4.775\times10^{-48}$ & $5.001\times10^{-41}$ \\
120 & 1 & $6.766\times10^{-39}$ & $2.23\times10^{-35}$ \\
120 & 2 & $3.961\times10^{-56}$ & $2.338\times10^{-52}$ \\
120 & 4 & $3.341\times10^{-93}$ & $3.469\times10^{-86}$ \\
% END GENERATED VERIFICATION TABLE
\bottomrule
\end{tabular}
\end{table}

For example, at $q=2,n=120$ the ratio of the left side of
\eqref{eq:firstequiv} to its proposed leading term is approximately
$1.0000561453$ for necklaces and $1.0000280316$ for primitive necklaces.
At $q=3,n=210$ the deviations of these ratios from one are approximately
$3.14\times10^{-15}$ and $1.57\times10^{-15}$, respectively.
These observations agree with the proved asymptotic estimate. At smaller
labels the ratio can be substantially different from one; the bound and
proof, not a visual numerical trend, justify the asymptotic statement.

The audits are finite tests, not substitutes for the general proofs.
In particular they do not establish the Banach completeness theorem,
the infinite support argument, or the uniform asymptotic statements.
No Lean kernel check of the new theorems was performed.

\section{Consequences and interpretation}\label{sec:consequences}

\subsection{A usable answer at an infinite action boundary}

The construction supplies two complementary kinds of finiteness.
For algebraic equality at one action, the support lemma makes the
coefficient extraction finite. For evaluation across an entire
accumulating block, the action--label norm supplies absolute convergence.
Neither mechanism can replace the other: a well-ordered support does
not by itself make its evaluation summable, and pointwise convergence
on one sheet does not by itself provide a faithful coefficient calculus.

The inverse theorem then uses a third organization, homogeneous degree
in the primitive input. Its analytic majorant remains geometric even
though a bounded-action prefix of the universal input is infinite.
This separation of exact action extraction, arithmetic labeling, and
analytic homogeneous summation is the main structural contribution.

\subsection{A precise beyond-all-orders effect}

The leading range displacement $-E_n(X)/p$ already contains infinitely
many universal action sectors. Once that entire block is retained,
\eqref{eq:firstcut} detects a specific smaller effect at its accumulation
boundary. Its coefficient is neither a freely chosen exponentially
small constant nor an interpolation-dependent Stokes parameter. It is
forced by the exact inverse equation and the identity $e_2^2=e_2$.
The factor $X^{-1}$ is forced by the slowly varying part of the core.

More generally, \eqref{eq:bresonance} separates two possible reasons an
action coefficient can vanish. The universal analytic kernel may vanish
in the pure-exponential limit, or the arithmetic sum over its lcm-labeled
tuples may cancel. Curvature resolves the first mechanism at order
$X^{-1}$ when $b\ne0$; it does not automatically resolve the second.
This distinction remains meaningful at action collisions such as those
in \eqref{eq:phicollisions} and \eqref{eq:mucollisions}.

\subsection{What remains outside the theorem}

The article does not construct the largest possible arithmetic
transseries ring, prove a general closure theorem for arbitrary
compositions, or define divisibility for noninteger labels. The
norm criterion is exact for the stated elementary block, not a
classification of every conditional or regularized summation method.
The inverse theorem is local in the analytic coordinate and is
uniform only in its stated large-core domains. A critical point near
$X=1/a$ is excluded. The curvature formulas are fixed-order
expansions, and the numerical checks use no interval arithmetic.
These boundaries leave several substantive research problems.

\section{Further research questions}\label{sec:questions}

\begin{question}[Joint action--label admissibility beyond one block]
For a general array $c_{s,l}$, characterize those formal arithmetic Hahn
series that admit a faithful holomorphic sheet realization with
uniform nonlinear inversion. Can an intrinsic condition, independent of
a chosen family $h_N(l)$, recover the useful part of the weighted
completion without forcing absolute $\ell^1$ summability?
\end{question}
A useful first benchmark is to compare the present norm with a norm
formed after summing absolute coefficients on each divisor fiber.
The latter can be substantially smaller, but its multiplicative and
faithfulness properties require separate analysis. Conditional
summation should not be conflated with the elementary-block criterion
already proved.

\begin{question}[Difference operations on arithmetic labels]
Construct a compatible arithmetic difference calculus generated by
congruence indicators $e_{d,r}(n)=\ind_{n\equiv r\pmod d}$, in which
integer shifts act naturally. Determine which joint action--congruence
weights make multiplication and shift uniformly continuous.
\end{question}
Products of congruence indicators are either zero or another congruence
indicator, according to the Chinese remainder compatibility condition.
This suggests an extension of the lcm rule, but Proposition
\ref{prop:idempotent} still prevents identifying label shifts with the
analytic exponential of a derivation. A successful theory should expose
both operations rather than hide this distinction.

\begin{question}[Uniform large homogeneous order]
Determine the joint asymptotics of $C_m(ka,X)$ when $m$ and $k$ grow
with $X$. At what scale does the first curvature term in
\eqref{eq:resonant} cease to dominate, and what rescaled profile replaces
it? The exact representation \eqref{eq:transformedC} supplies a concrete
starting point for a two-parameter saddle analysis.
\end{question}
The present theorem controls fixed $m,k$ only. Expanding in $1/X$ and
then allowing $m$ to grow is not justified by its proof. A uniform
result could connect geometric homogeneous convergence to sharp
optimal truncation within individual action families.

\begin{question}[Arithmetic cancellation at resonant actions]
Classify the labels and coefficient sequences for which
\[
 \sum_{\lambda_{d_1}+\cdots+\lambda_{d_m}=ka}
 \left(\prod_i c_{d_i}\right)e_{\lcm(d_1,\ldots,d_m)}
\]
vanishes identically, or on specified arithmetic subsequences. How do
totient positivity and M\"obius signs alter the first surviving
curvature-lifted term?
\end{question}
The finite action algorithm makes each fixed case decidable. A
structural classification would require more than checking initial
coefficients, especially where several lcm labels survive at the same
action.

\begin{question}[More complicated accumulation sets]
Extend the action--label completion to positive generator sets with
several finite accumulation levels or higher order types. Identify
conditions that separately guarantee well-ordering, finite exact-action
decompositions, and convergence of the label-weighted norm.
\end{question}
The positive lower bound $a/2$ bounds homogeneous length in the present
proof. Removing it can destroy this argument even when individual
labels leave only finitely many input terms. A replacement theorem
would need to control both arbitrarily long products and any action
accumulation at zero.

\begin{question}[Optimal expanding domains]
For $c_d$ close to the threshold $|c_d|\asymp\ee^{ad}$, determine the
smallest label-dependent analytic half-planes on which inverse families
remain uniformly summable. Can the norm be sharpened sufficiently to
describe logarithmic or iterated-logarithmic boundary corrections?
\end{question}
For example, elementary inputs with
$|c_d|=\ee^{ad}/d^\gamma$ are admitted exactly for $\gamma>1$ in the
present norm. At or below that boundary, every fixed sheet still exists;
the issue is the geometry and summability of the family, not the
existence of the finite divisor sum.

\begin{question}[Certified numerical realization]
Implement the jet algorithm with outward-rounded intervals, including
the large Lambert core and an explicit divisor-coefficient majorant.
Can the majorant be sharpened by action-aware estimates to provide
near-optimal certified errors without enumerating all tuples?
\end{question}
The current Cauchy bound is explicit but can be conservative, as
Table~\ref{tab:errors} shows. A proof-producing implementation should
separate the exact combinatorial count, interval evaluation of the
core, and propagation of the inverse tail.

\begin{question}[Formal verification of the arithmetic completion]
Formalize the action-support theorem, lcm-coordinate ring, weight
inequality, and homogeneous coefficient identities in Lean. Then
connect a finite checker to the mathematical semantics of the jet
algorithm and to the analytic tail theorem.
\end{question}
Appendix~\ref{app:lean} gives a dependency order and the limited existing
repository interfaces that can be reused. The purpose is to verify the
new construction, not to relabel an unverified analytic argument as an
already formalized theorem.

\begin{question}[Critical cores with arithmetic sector accumulation]
Study inverse sheets when the core approaches its critical slope
$a-1/X=0$ while an arithmetic primitive block remains summable. Is there
a uniform fold coordinate compatible with the full lcm-labeled block,
and how does it interact with the action accumulation cut?
\end{question}
The geometric disk proof here deliberately assumes a lower slope bound.
Existing repository fold analyses concern other models. Combining such
coordinates with the present arithmetic summation would be a separate
extension, not a consequence of the large-$X$ theorem.

\section{Conclusion}

An accumulating action support need not prevent exact symbolic
multiplication, and arithmetic indicators need not be replaced by
trigonometric surrogates to obtain a summable analytic theory. The
lcm-coordinate ring, a well-ordered positive action monoid, and a weight
that couples action to the earliest possible arithmetic label together
provide a workable alternative. The resulting inverse expansion is
absolutely convergent on specified sheets and exact on the counting
sequence's range.

The first post-cut inverse sector illustrates why this framework is
more than a reformulation of finite divisor cutoffs: the whole
primitive block must first be retained, after which a curvature-induced
$X^{-1}\ee^{-aX}$ term becomes visible. The all-order kernel identity
explains its resonance mechanism and its extension to a general
power-law core. These are precise research results within an explicit
scope; their broader novelty, optimized numerical realization, and
formal verification remain tasks for further work.

\appendix
\section{Repository provenance and scope ledger}\label{app:provenance}

The repository was inspected through its GitHub connector, with a
subsequent pinned read of the exact open-question passage. The snapshot
used for this article is
\begin{center}
\small\snapshot.
\end{center}
Its parent was
\texttt{04473354a0f3edff2d6c365caf26170ca6b88735}.
The snapshot's final commit concerns unrelated ordinal reports; the
transseries source blob recorded below is unchanged by that commit.
The article does not depend on any unrelated claims in the repository.

\paragraph{Primary mathematical source.}
{\small\path{Analysis/Transseries/docs/series-and-transseries/Combinatorial_Transseries_Inverses/Combinatorial_Transseries_Inverses.tex}}\par
The source blob is
\texttt{17647298fa31d61c061d556a3b1e451a228ae20e}.
The necklace chapter and its finite cutoffs were read, along with the
synthesis and research directions. The explicit arithmetic-algebra
question was reread at the pinned snapshot, source lines 4955--4978.

\paragraph{Comparison with existing arrivals.}
{\small\path{Analysis/Transseries/docs/series-and-transseries/README.md}}\par
Its entry for
\path{Action_Accumulation_Nonlinear_Inversion} states that that
article does not retain divisibility indicators. The catalogue was used
to avoid duplicating an analytic action-measure result. No complete
audit of all repository research reports was attempted.

\paragraph{Existing Lean source inspected.}
{\small\path{Analysis/Transseries/Lean/Transseries/TransseriesWellBased.lean}}\par
Its blob is
\texttt{06d5ea1663573d96d5712636347d44764b40f800}.
The actual declarations, rather than only a catalogue description,
were read for the formalization comparison below.


The package also includes \texttt{source\_notes/provenance.json} and
\texttt{source\_notes/scope.md}. These record what was consulted and what
was not claimed. No copy of the entire repository or its large source
volumes is required to compile this article.

\section{Lean interfaces and a proposed verification route}\label{app:lean}

The inspected module places its declarations in namespace
\texttt{Fabius}, although its module path is now under
\texttt{Transseries}. It supplies
\path{dickson_isPWO},
\path{dickson_antichain_finite}, and
\path{dickson_isPWO_pi}; the Neumann interfaces include
\path{neumann_isPWO} and
\path{neumann_finite_factorizations}, together with additive
versions generated by \texttt{to\_additive} and explicit order-dual
wrappers \cite{repolean}.

The source describes these as bridges to Mathlib's partially
well-ordered-set and finite-antidiagonal results. Their presence is
useful, but it is important not to overstate it. They do not instantiate
the lcm-labeled Banach algebra, prove its norm inequality, verify the
new analytic inverse theorem, or establish the curvature formula.
No such formalization is delivered with this article.

A proposed dependency order is as follows.
\begin{enumerate}
\item Define the rational normalized actions $1-1/d$ for $d\ge2$.
Use the finite-length bound and the available Dickson interface to
prove well-ordering and finite exact-action representations. Keep the
embedding into real actions $a\alpha$ separate from the rational
combinatorial support.
\item Define divisor-coordinate arithmetic functions and prove
$e_i e_j=e_{\lcm(i,j)}$. Establish the finite lcm convolution, its
associativity, and the divisor-zeta equivalence. This stage needs no
analysis or asymptotics.
\item Construct the coefficient ring on the action monoid, proving
finiteness before defining each convolution. Verify the universal
polynomial identities underlying finite Lagrange jets over a
commutative rational algebra.
\item Introduce the weighted summable-family space. Prove the single
weight inequality from Theorem~\ref{thm:banach}, then use it to build
the normed algebra and its evaluation homomorphisms. Establish the
admission equivalence by comparison of positive series.
\item Formalize the analytic disk bounds, root existence and uniqueness,
and Cauchy coefficient estimates. Finally connect the finite jet
checker to this analytic theorem. The finite program should not be
trusted as a source of unproved universal identities.
\end{enumerate}

The exact derivative-of-binomial identity in
Theorem~\ref{thm:resonance} and the finite action search are plausible
early verification targets. The uniform holomorphic-family construction
is a substantially larger step. This is an implementation roadmap,
not a report of a Lean build.

\section{Reproducibility and deliverable contents}\label{app:repro}

The source is a standalone LaTeX article. From its directory, run
\begin{verbatim}
python verification/verify.py
pdflatex -interaction=nonstopmode -halt-on-error arithmetic_transseries.tex
pdflatex -interaction=nonstopmode -halt-on-error arithmetic_transseries.tex
pdflatex -interaction=nonstopmode -halt-on-error arithmetic_transseries.tex
\end{verbatim}
The equivalent convenience command is \texttt{sh build.sh}. A recent
Python with \texttt{mpmath} and a standard TeX installation with the
packages used in the preamble are required. The generated verification table is embedded in this standalone source;
a matching table fragment is also shipped. The article can therefore be
compiled without rerunning the numerical script.

The verification program records its working precision and separates
exact checks from numerical diagnostics. The counts in this article
refer to the program included in this package, not to any pre-existing
repository audit. The PDF was compiled from the delivered source and
its rendered pages were inspected. No font files, third-party source
volumes, or build-cache files are included in the archive.

\begin{thebibliography}{9}

\bibitem{repo}
Vladimir Reshetnikov, \emph{ProveIt: Combinatorial Transseries and Their
Inverses}, repository manuscript, pinned snapshot of 29 September 2026.
Primary path, source blob, and inspected passages are specified in
Appendix~\ref{app:provenance}.
\url{https://github.com/VladimirReshetnikov/ProveIt/tree/db68f0853c3c69cd930caedab6f6ad9addb11eaf/Analysis/Transseries}.

\bibitem{repoindex}
Vladimir Reshetnikov, \emph{Series and transseries}, ProveIt catalogue,
including the arrivals of 29 September 2026 and the scope description of
\emph{Action Accumulation and Nonlinear Inversion}. Repository and
snapshot as in \cite{repo}; path
\path{Analysis/Transseries/docs/series-and-transseries/README.md}.

\bibitem{repolean}
Vladimir Reshetnikov, \emph{TransseriesWellBased.lean}, ProveIt Lean
source, snapshot as in \cite{repo}; path and blob in
Appendix~\ref{app:provenance}.

\bibitem{toth}
L\'aszl\'o T\'oth, Multiplicative arithmetic functions of several
variables: a survey, in \emph{Mathematics Without Boundaries: Surveys in
Pure Mathematics}, T.~M. Rassias and P.~M. Pardalos (eds.), Springer,
2014, pp.~483--514. In particular, Section 4.4 on lcm convolution.
\url{https://arxiv.org/abs/1310.7053}.

\bibitem{edgar}
G.~A. Edgar, Transseries for beginners, \emph{Real Analysis Exchange}
\textbf{35} (2009/2010), 253--310.
\url{https://arxiv.org/abs/0801.4877}.

\bibitem{chebolu}
Sunil K. Chebolu, J\'an Min\'a\v{c}, Tung T. Nguyen, and Nguyen Duy Tan,
\emph{Analytic Properties of Necklace Polynomials}, arXiv preprint,
version 1, 12 May 2026.
\url{https://arxiv.org/abs/2605.11445}.

\end{thebibliography}
\end{document}
